\documentclass[11pt]{article}
\usepackage[margin=1in]{geometry}
\usepackage{amsmath,amssymb,amsthm}
\usepackage{physics}
\usepackage{enumitem}
\usepackage{hyperref}
\usepackage{titlesec}

\theoremstyle{definition}
\newtheorem{definition}{Definition}[section]
\newtheorem{example}{Example}[section]
\newtheorem{remark}{Remark}[section]
\newtheorem{proposition}{Proposition}[section]
\newtheorem{theorem}{Theorem}[section]

\titleformat{\section}{\large\bfseries}{Lecture 8.\thesection}{1em}{}

\title{\textbf{Lecture 8: Grover's Algorithm} \\ \large Understanding Quantum Information \& Computation}
\author{Based on the lectures by John Watrous (IBM Quantum)}
\date{}

\begin{document}
\maketitle
\tableofcontents
\newpage

%----------------------------------------------------------------------
\section{Overview}
%----------------------------------------------------------------------

This lecture covers \textbf{Grover's algorithm}, which provides a \emph{quadratic} speedup for unstructured search problems.  Topics:

\begin{enumerate}
    \item The unstructured search problem.
    \item The Grover iteration and its geometric interpretation.
    \item Optimal number of iterations (unique and multiple solutions).
    \item Lower bounds and remarks.
\end{enumerate}

%----------------------------------------------------------------------
\section{The Unstructured Search Problem}
%----------------------------------------------------------------------

\begin{definition}
Given a function $f:\{0,1\}^n\to\{0,1\}$, find a \textbf{solution}: a string $z$ with $f(z)=1$, or report that none exists.
\end{definition}

Let $N=2^n$ be the total number of input strings.

\begin{itemize}
    \item The function $f$ may be given as a black box (query model) or as a Boolean circuit.
    \item No structure is assumed: it is like searching an unsorted phone book by number.
    \item This is an NP-complete problem when $f$ is specified by a Boolean circuit.
\end{itemize}

\paragraph{Classical cost.}  Deterministic: $N$ queries (worst case).  Randomised: still $\Theta(N)$.

\paragraph{Quantum cost (Grover).}  $O(\sqrt{N})$ queries---a \textbf{quadratic advantage}.

\subsection{Phase Query Gates}

\begin{definition}
The \textbf{phase query gate} for $f$ is
\[
    Z_f\ket{x} = (-1)^{f(x)}\ket{x}.
\]
It can be implemented from a standard query gate $U_f$ via phase kickback (using a $\ket{-}$ ancilla).
\end{definition}

We also define the phase gate for the OR function on $n$ bits:
\[
    Z_{\mathrm{OR}}\ket{x} = \begin{cases}-\ket{x} & x\neq 0^n,\\ \;\;\ket{x} & x=0^n.\end{cases}
\]

%----------------------------------------------------------------------
\section{Grover's Algorithm}
%----------------------------------------------------------------------

\subsection{The Algorithm}

\begin{enumerate}
    \item \textbf{Initialise:} $\ket{\psi_0}=H^{\otimes n}\ket{0^n}=\frac{1}{\sqrt{N}}\sum_{x\in\{0,1\}^n}\ket{x}$.
    \item \textbf{Iterate:} Apply the \textbf{Grover operator} $G$ a total of $t$ times.
    \item \textbf{Measure} in the standard basis to obtain a candidate solution $x$.
\end{enumerate}

\subsection{The Grover Operator}

\[
    G = -H^{\otimes n}\,Z_{\mathrm{OR}}\,H^{\otimes n}\,Z_f.
\]
Equivalently, $G = (2\ket{\psi_0}\!\bra{\psi_0}-I)\,Z_f$, where:
\begin{itemize}
    \item $Z_f$ is the \textbf{oracle}: flips the phase of solutions.
    \item $2\ket{\psi_0}\!\bra{\psi_0}-I$ is the \textbf{diffusion operator}: reflects about the uniform superposition.
\end{itemize}

Each application of $G$ requires \textbf{one query} to $f$.

%----------------------------------------------------------------------
\section{Geometric Analysis}
%----------------------------------------------------------------------

\subsection{Two-Dimensional Invariant Subspace}

Partition $\{0,1\}^n$ into solutions $A_1=\{x:f(x)=1\}$ and non-solutions $A_0=\{x:f(x)=0\}$.  Let $|A_1|=M$ (number of solutions).  Define normalised states:
\[
    \ket{A_0} = \frac{1}{\sqrt{N-M}}\sum_{x\in A_0}\ket{x},
    \qquad
    \ket{A_1} = \frac{1}{\sqrt{M}}\sum_{x\in A_1}\ket{x}.
\]

The initial state lies in the plane spanned by $\ket{A_0}$ and $\ket{A_1}$:
\[
    \ket{\psi_0} = \cos\frac{\alpha}{2}\,\ket{A_0}+\sin\frac{\alpha}{2}\,\ket{A_1},
    \qquad\text{where}\quad \sin\frac{\alpha}{2}=\sqrt{\frac{M}{N}}.
\]

\subsection{Grover Operator as a Rotation}

\begin{theorem}
Restricted to $\operatorname{span}\{\ket{A_0},\ket{A_1}\}$, the Grover operator $G$ is a \textbf{rotation by angle $\alpha$}:
\[
    G^t\ket{\psi_0} = \cos\!\left(\frac{(2t+1)\alpha}{2}\right)\ket{A_0}
                     +\sin\!\left(\frac{(2t+1)\alpha}{2}\right)\ket{A_1}.
\]
\end{theorem}

The probability of measuring a solution after $t$ iterations is
\[
    p(t) = \sin^2\!\left(\frac{(2t+1)\alpha}{2}\right).
\]

%----------------------------------------------------------------------
\section{Choosing the Number of Iterations}
%----------------------------------------------------------------------

\subsection{Unique Solution ($M=1$)}

Here $\sin(\alpha/2)=1/\sqrt{N}$, so $\alpha\approx 2/\sqrt{N}$ for large $N$.

The success probability is maximised when $(2t+1)\alpha/2\approx\pi/2$, giving
\[
    t^* = \left\lfloor\frac{\pi}{2\alpha}-\frac{1}{2}\right\rfloor
        \approx \frac{\pi}{4}\sqrt{N}.
\]

\begin{theorem}
Grover's algorithm finds the unique solution with probability $\ge 1-1/N$ using $O(\sqrt{N})$ queries.
\end{theorem}

\subsection{Multiple Solutions ($M>1$)}

With $M$ solutions, $\alpha = 2\arcsin\sqrt{M/N}$ and the optimal iteration count is
\[
    t^* \approx \frac{\pi}{4}\sqrt{\frac{N}{M}}.
\]

If $M$ is known, we can choose $t^*$ directly.  If $M$ is unknown, \textbf{exponential search} strategies can be employed: try geometrically increasing values of $t$ until a solution is found.

\begin{remark}[Over-rotation]
Applying $G$ \emph{too many} times causes the state to rotate \emph{past} $\ket{A_1}$ and the success probability decreases.  Grover's algorithm is not monotone---more iterations can be worse.
\end{remark}

%----------------------------------------------------------------------
\section{Optimality}
%----------------------------------------------------------------------

\begin{theorem}[BBBV lower bound]
Any quantum algorithm that solves the unstructured search problem with bounded error must make $\Omega(\sqrt{N})$ queries.
\end{theorem}

Hence Grover's algorithm is \textbf{optimal} (up to constant factors) in the query model.

%----------------------------------------------------------------------
\section{Remarks and Generalisations}
%----------------------------------------------------------------------

\begin{itemize}
    \item \textbf{Amplitude amplification.}  Grover's technique generalises: given any quantum algorithm that succeeds with probability $p$, amplitude amplification boosts the success probability to near~$1$ using $O(1/\sqrt{p})$ repetitions.
    \item \textbf{Practical relevance.}  A quadratic speedup is significant but modest compared to the exponential speedup of Shor's algorithm.  Whether Grover-type speedups will be practically useful depends on hardware advances.
    \item \textbf{No structure exploited.}  Grover's algorithm treats $f$ as a black box.  For structured problems, more powerful algorithms (like Shor's) may apply.
    \item \textbf{Quantum minimum finding.}  Grover search can be adapted to find the minimum of a function, yielding a quadratic speedup for optimisation problems.
\end{itemize}

%----------------------------------------------------------------------
\section{Summary}
%----------------------------------------------------------------------

\begin{center}
\begin{tabular}{lcc}
\hline
& \textbf{Classical} & \textbf{Quantum (Grover)} \\\hline
Queries for search ($M=1$) & $\Theta(N)$ & $\Theta(\sqrt{N})$ \\
Queries for search ($M$ solutions) & $\Theta(N/M)$ & $\Theta(\sqrt{N/M})$ \\
\hline
\end{tabular}
\end{center}

Key ideas:
\begin{itemize}
    \item The Grover operator is a \textbf{rotation} in a 2D subspace.
    \item Repeated application steers the state toward solutions.
    \item The $\sqrt{N}$ query count is \textbf{provably optimal} for unstructured search.
\end{itemize}

\end{document}
